\paragraph{AQ-136-DERIVED-COMPATIBILITY: a uniform necessary invariant for simultaneous completion.}
Let $A$ be any simplicial bigraded strict PD differential
$S_k$-algebra; no flatness or admissibility condition on $A$
is imposed.  Let $C_1,\ldots,C_s\in\mathcal A_k$ and fix
strict maps $f_i:A\to C_i$.  Suppose these maps have a
simultaneous admissible completion:
\[
 A\xrightarrow{g}R\xrightarrow{j_i}C_i,
       \qquad R\in\mathcal A_k,\qquad j_i g=f_i.
                                                               \tag{136.DC1}
\]
There is a necessary condition imposing compatibility
after derived coefficient change, in addition to the
integral comparison.
For every $d\ge0$ there is a weight bound $W_d$ such that,
simultaneously for all $w\ge W_d$, all finite simplicial
sets $X$, all integers $n\ge2$, and all $0\le q\le d$,
the maps
\[
 H^q\bigl(A^{w,*}(X)\otimes^{\mathbf L}_{\mathbb Z}
                         \mathbb Z/n\bigr)
 \longrightarrow \mathcal T^q_n(X;k)\otimes_k S^w
                                                               \tag{136.DC2}
\]
induced by the $f_i$ and the canonical admissible comparisons
are equal for every $i$.  Here the target has the following
precise, theory-independent meaning.  Put
\[
 \mathcal E_X(k)=R\operatorname{Hom}_{s\operatorname{Mod}k}
                         (k[X],\underline{k}),\qquad
 \mathcal T^q_n(X;k)=H^q\bigl(\mathcal E_X(k)
                  \otimes^{\mathbf L}_{\mathbb Z}\mathbb Z/n\bigr).
                                                               \tag{136.DC3}
\]
The underline denotes the constant simplicial module.  In
particular these targets have natural exact sequences
\[
 0\longrightarrow H^q(X;k)/n
 \longrightarrow\mathcal T^q_n(X;k)
 \longrightarrow H^{q+1}(X;k)[n]\longrightarrow0.
                                                               \tag{136.DC4}
\]
No splitting of this sequence is chosen.  In (136.DC2), the
notation with $S^w$ uses its specified free coefficient
generator $T^{[w]}$; equivalently one may replace
$\underline{k}$ in (136.DC3) by $\underline{S^w}$.

The bound $W_d$ is independent of $n$ and $X$.
It can depend on the proposed completion $R$ and on the
finite target family.  Consequently a discrepancy for
arbitrarily large weights, even for one fixed $X$, one
fixed prime modulus and one fixed degree, rules out that
simultaneous completion.  This is a necessary condition,
not a sufficiency criterion for original common sources.

\begin{proof}
We first specify the derived coefficient operation on a
possibly nonflat complex.  For any cochain complex $M$ of
abelian groups define
\[
 D_n(M)^q=M^q\oplus M^{q+1},\qquad
 d(a,b)=(da+n b,-db).
                                                               \tag{136.DC5}
\]
This is the tensor product with the free resolution
$[\mathbb Z\xrightarrow{n}\mathbb Z]$ of $\mathbb Z/n$,
placed in degrees $-1,0$, using the equivalent mapping-cone
sign convention.  A bounded complex of free abelian groups
is $K$-flat: tensoring an acyclic complex with each free
term is exact, and the finite total complex is acyclic by
the finite filtration.  Thus (136.DC5) computes
$M\otimes^{\mathbf L}_{\mathbb Z}\mathbb Z/n$ for every
$M$, with no flatness assumption on its terms.
The mapping-cone long exact sequence gives, naturally,
\[
 0\longrightarrow H^q(M)/n
 \longrightarrow H^q(D_n(M))
 \longrightarrow H^{q+1}(M)[n]\longrightarrow0.
                                                               \tag{136.DC6}
\]
It follows that a map which is an isomorphism on integral
cohomology in degrees $q,q+1$ is an isomorphism on
$H^q(D_n(-))$, for every $n$, with no change to its
cohomological range.  This follows directly from the
commutative diagram of (136.DC6), by the short five lemma.

We use the actual natural comparison construction in
AQ-136-COMPARISON, rather than choosing unrelated
isomorphisms of cohomology groups.  Write
$\Gamma_X(F)=\operatorname{Hom}_{s\operatorname{Mod}k}(k[X],F)$.
For any $C\in\mathcal A_k$ there are natural maps in the
derived category of $k$-modules
\[
 C^{w,*}(X)=\Gamma_X(C^{w,*})
       \longrightarrow R\Gamma_X(C^{w,*})
       \longleftarrow R\Gamma_X(\underline{S^w}).
                                                               \tag{136.DC7}
\]
The right-hand map is induced by the coefficient
augmentation $\underline{S^w}\to C^{w,*}$, and the
left-hand map is the ordinary-to-derived-section map.
Both are therefore natural for every strict
$S$-preserving morphism $C\to C'$.
The complete proof of AQ-136-COMPARISON gives one weight
cutoff, independent of $X$, at which both arrows are
isomorphisms on cohomology through any prescribed finite
degree.  Indeed its first hypercohomology spectral sequence
uses the uniform vertical cohomology bounds, and its
second uses the acyclicity of each required horizontal
row for all higher derived sections.  All total degrees
in that argument have finite diagonals.  Thus this is a
statement about the two displayed maps themselves, not
only the existence of abstract cohomology isomorphisms.
The same proof identifies
$H^q(R\Gamma_X(\underline{k}))$ naturally with $H^q(X;k)$.

For the finite family $R,C_1,\ldots,C_s$, take the maximum
of their comparison cutoffs through degree $d+1$.
Apply $D_n$ to the two arrows in (136.DC7).
By (136.DC6), their cohomology maps through degree $d$
are isomorphisms for every $n\ge2$.  Inverting these
cohomology isomorphisms gives the canonical comparisons
with (136.DC3), and (136.DC6) applied to the coefficient
object proves (136.DC4).
Since every $j_i$ preserves the coefficient map, naturality
of (136.DC7) says that it acts as the identity on the
common target (136.DC3), including after $D_n$.
After precomposition with $D_n(g)$, all the maps in
(136.DC2) are therefore equal.  The weight cutoff was
chosen before $n$ and $X$, so it has the asserted uniformity.

This argument performs derived coefficient change
\emph{after} taking sections.  It never replaces
$A(X)\otimes^{\mathbf L}\mathbb Z/n$ by ordinary reduction
$A(X)/n$, nor identifies it with sections of a
coefficient-changed simplicial algebra.  Such an
interchange could fail for a nonflat source.  In fact
the comparison proof used above works for arbitrary
simplicial sets, so the stated finite-$X$ version needs
no additional finiteness or base-change theorem.
\end{proof}

\paragraph{The earlier obstruction is exactly a failed derived compatibility test.}
Retain the explicit $A$, $f_0,f_1:A\to B$ from
AQ-136-HIGHER-COMPLETION, over $k=\mathbb Z$ and at a
fixed prime $p$.  On the simplicial circle $X$, its
elements $v_w,u_w$ satisfy
$dv_w=p u_w$ and $du_w=0$.  Hence
\[
                 (v_w,-u_w)\in D_p(A^{w,*}(X))^1
                                                               \tag{136.DC8}
\]
is a closed element by (136.DC5).  Its image under $f_0$
is zero.  Under $f_1$ it is
$(T^{[w-1]}dx,0)$, whose class is the image of
$wT^{[w]}$ under the left-hand injection of (136.DC6).
The integral endpoint calculation (136.HC8) gives this
normalization without any coefficient-change assumption.
Since $H^1(B^{w,*}(X))=\mathbb ZT^{[w]}$ and
$H^2(B^{w,*}(X))=0$ in the comparison range, this class
is nonzero exactly when $p\nmid w$.
Such weights are arbitrarily large.  Thus the uniform
necessary condition (136.DC2) already contradicts the
existence of (136.HC2).  This recovers the checked
obstruction with a single mod-$p$ cone class, without
an assertion that the source is flat or that ordinary
mod-$p$ sections compute its derived coefficient change.
For a proposed positive source construction, retaining
only its integral $H^0/H^1$ and rational comparison is
therefore insufficient: the prescribed maps must also
preserve these derived coefficient compatibilities.
